\chapter{Sensors and Microsoft.Devices}
\label{ch:sensors}
\label{ch:devices-sensors}

\cnans{Microsoft::Devices::Sensors} ports the Windows Phone 7 sensor model to SDL3 and Android
NDK inputs. The CNA-original host utility layer in Chapter~\ref{ch:host-devices} is separate.
Sensor callbacks may run outside the game thread, which makes lifecycle and synchronization part
of the public usage contract.

\section{Sensor model and delivery}

\cnaclass{SensorBase<TSensorReading>} supplies \texttt{CurrentValue}, \texttt{IsDataValid},
\texttt{TimeBetweenUpdates} (default 2 ms), \texttt{CurrentValueChanged}, and abstract
\texttt{Start}/\texttt{Stop}. Do not touch \cnaclass{GraphicsDevice} or other main-thread-only
objects in a sensor handler; transfer the reading through a synchronized queue.

\begin{lstlisting}
Accelerometer accelerometer;
accelerometer.CurrentValueChanged.Add(
    [](System::Object*,
       const SensorReadingEventArgs<AccelerometerReading>& args) {
        Vector3 acceleration =
            args.getSensorReadingProperty().getAccelerationProperty();
        // Enqueue acceleration for the game thread.
    });
accelerometer.Start();
// ... later ...
accelerometer.Stop();
accelerometer.Dispose();
\end{lstlisting}

Desktop Accelerometer and Gyroscope install an SDL event watch rather than a polling thread.
\texttt{SDL\_EVENT\_SENSOR\_UPDATE} is handled on the thread that pushes it. Acceleration is
converted from SI units to $g$ by division by 9.80665; angular velocity is already in radians
per second. Readings receive UTC wall-clock timestamps.

Android Compass and Motion use \texttt{ASensorManager}, event queues, and \texttt{ALooper} on
worker threads, without JNI. A Motion instance can own six workers for rotation, game rotation,
gravity, linear acceleration, gyroscope, and magnetic field. The native route expects Android
API 24 or later.

\begin{center}
\begin{tabularx}{\textwidth}{@{}lXX@{}}
\toprule
\textbf{Type} & \textbf{Implemented route} & \textbf{Boundary} \\
\midrule
\cnaclass{Accelerometer} & SDL3 on desktop, Android, and iOS & WP7 legacy
\texttt{ReadingChanged} also fires \\
\cnaclass{Gyroscope} & SDL3 sensor events & CNAEXT \texttt{State}; no legacy event \\
\cnaclass{Compass} & Android NDK only & Unsupported elsewhere; \texttt{Start()} throws \\
\cnaclass{Motion} & Android NDK fusion only & Unsupported elsewhere; six possible workers \\
\bottomrule
\end{tabularx}
\end{center}

\section{Concurrency contract}

Disposal is claimed atomically so only one racing caller performs cleanup. A losing caller
waits for terminal state, and an RAII guard publishes that state even if cleanup throws.
Accelerometer and Gyroscope track callbacks in flight so a handler can dispose its own sensor
without waiting for itself.

Current value and validity are published under one lock by
\texttt{SetCurrentValueAndMarkDataValid}. Separate locks previously allowed a reader to observe
\texttt{IsDataValid=true} with the prior value. Update throttling compares elapsed time at the
configured interval's coarser resolution; the earlier reverse conversion overflowed signed
64-bit ticks for \texttt{TimeSpan.MaxValue}.

Compass and Motion use a two-phase start/stop protocol: reserve under the owner mutex, invoke
the backend without that lock, then commit or roll back. A condition variable prevents a new
start while an orphaned start's cleanup call is still in flight. Stop deliberately remains able
to supersede a start without waiting on it.

\begin{historynote}{Why sanitizer evidence matters here}
The first lifecycle rewrite appeared race-free under ordinary tests and manual review.
ThreadSanitizer found that Stop cleared a transition flag before an orphaned Start had finished
its backend cleanup, allowing a third Start to enter concurrently. The stress test reproduced
both the race and a use-after-free in callback bookkeeping. The in-flight-call counter and
condition variable are the current fix. UndefinedBehaviorSanitizer separately found the
update-interval overflow.
\end{historynote}

Exceptions cannot cross SDL, Android C callbacks, or raw thread entry points. CNA records
backend, operation, native error, device, timestamp, and severity through a \texttt{noexcept}
diagnostic channel. The Android bridge distinguishes \texttt{std::exception::what()} from an
unknown throw. Other device paths still retain older ad hoc counters, so the shared diagnostic
record is not yet universal.

\section{Compass and motion math}

Compass uses upright heading when gravity meets the $\cos 45^\circ$ threshold; otherwise it
derives flat heading from the quaternion. Results are normalized to $[0,360)$, and an invalid
quaternion maps to $0^\circ$. Accuracy values High/Medium/Low/Unreliable map to 5/15/20/180
degrees. Without declination data, \texttt{TrueHeading} equals \texttt{MagneticHeading}.

Motion normalizes Android's rotation quaternion and derives pitch, yaw, and roll. A process-wide
CNAEXT landscape remap affects acceleration, rotation rate, gravity, and linear acceleration,
but not \texttt{Attitude}: the requested axis reflection has determinant $-1$ and cannot be
represented by a quaternion. Fusion accepts the latest readings within 500 ms; it does not
interpolate them to one timestamp.

\section{Vibration}

\cnaclass{VibrateController} is a process-lifetime singleton. \texttt{Start(TimeSpan)} accepts
the WP7 range $[0,5\text{s}]$; \texttt{Start(TimeSpan,float)} and
\texttt{StartLeftRight} are CNAEXT intensity controls. SDL's Android haptic route blends motor
intensities; independent left/right controller rumble is available through \cnaclass{GamePad},
not this phone-style API.

\begin{lstlisting}
VibrateController* vibration = VibrateController::getDefaultProperty();
if (vibration->getIsSupportedProperty()) {
    vibration->StartLeftRight(
        1.0f, 0.3f, System::TimeSpan::FromSeconds(0.2));
}
// Start(duration, 0.0f) still installs a zero-strength effect.
vibration->Stop();
\end{lstlisting}

\texttt{DevicesShutdownCoordinator::Shutdown()} must precede the application's
\texttt{SDL\_Quit()}. It prevents the singleton's later static destructor from closing an SDL
device after SDL has already released native state. One global mutex serializes sensor and
haptic subsystem initialization and preserves the documented per-instance-before-global lock
order.

\section{Known limitation and evidence status}

All four concrete sensors expose \texttt{Dispose(bool)} publicly even though the base declares
it protected. External \texttt{Dispose(false)} marks the instance disposed without stopping it,
decrementing use counts, or releasing SDL ownership. This is an open API/lifetime defect; call
the parameterless public \texttt{Dispose()}.

Focused tests cover math, throttling, lifecycle, fake backends, reentrancy, and TSan stress.
They do not constitute hardware QA. The Android source notes that its sign and zero conventions
have not been checked on a physical device, and the Android demo is an older copy of the desktop
demo. Android behavior is source-proven and host-tested where fakes permit, not device-verified.
